\begin{table*}[h!]
\small  % You can also use \footnotesize, \scriptsize, or \tiny
\centering
\begin{tabularx}{\textwidth}{|X|}
\hline
\#\#\# Demo 1\\ Question: How can I compress a custom ONNX model using singular value decomposition?\\ \\ Candidate Documents\\ API\_1: Converter.convert\_model(input\_model\_path: str, output\_dir: str, target\_framework: Union{[}str, Framework{]}, target\_device\_name: Union{[}str, DeviceName{]}, target\_data\_type: Union{[}str, DataType{]} = DataType.FP16, target\_software\_version: Optional{[}Union{[}str, SoftwareVersion{]}{]} = None, input\_shape: Optional{[}InputShape{]} = None, dataset\_path: Optional{[}str{]} = None, wait\_until\_done: bool = True): Convert a model to the specified framework. Returns model conversion task dictionary.\\ API\_2: Compressor.recommendation\_compression(compression\_method: CompressionMethod, recommendation\_method: RecommendationMethod, recommendation\_ratio: float, input\_model\_path: str, output\_dir: str, input\_shapes: List{[}Dict{[}str, int{]}{]}, framework: Framework = Framework.PYTORCH, options: Options = Options(), dataset\_path: Optional{[}str{]} = None): Compress a recommendation-based model using the given compression and recommendation methods. Returns source model and compressed model information.\\ API\_3: Compressor.automatic\_compression(input\_model\_path: str, output\_dir: str, input\_shapes: List{[}Dict{[}str, int{]}{]}, framework: Framework = Framework.PYTORCH, compression\_ratio: float = 0.5): Compress a model automatically based on the given compression ratio.\\ \\ Judgment: A\\ \\ By using above documents, you can compress the model using API\_2 Compressor.recommendation\_compression().\\ As documents support clear, accurate, and engaging answers, query is answerable so your output should be A.\\ \\ \#\#\# Demo 2\\ Question: Show me an example of training a ResNet model for image classification.\\ \\ Candidate Documents\\ API\_1: Trainer.set\_model\_config(model\_name: str, img\_size: int, use\_pretrained: bool = True, load\_head: bool = False, path: Optional{[}str{]} = None, fx\_model\_path: Optional{[}str{]} = None, optimizer\_path: Optional{[}str{]} = None): Set the model configuration for the Trainer.\\ API\_2: Trainer.set\_logging\_config(project\_id: Optional{[}str{]} = None, output\_dir: str = './outputs', tensorboard: bool = True, csv: bool = False, image: bool = True, stdout: bool = True, save\_optimizer\_state: bool = True, validation\_epoch: int = 10, save\_checkpoint\_epoch: Optional{[}int{]} = None): Set the logging configuration.\\ API\_3: Trainer.set\_training\_config(optimizer, scheduler, epochs: int = 3, batch\_size: int = 8): Set the training configuration.\\ \\ Judgment: P\\ \\ You can set ResNet model using API\_1 Trainer.set\_model\_config(), but you cannot select the task with above documents.\\ You can partially answer to the query, so your output should be P.\\ \\ \#\#\# Demo 3\\ Question: How do I set random seed for environmental setup?\\ \\ Candidate Documents\\ API\_1: Trainer.set\_logging\_config(project\_id: Optional{[}str{]} = None, output\_dir: str = './outputs', tensorboard: bool = True, csv: bool = False, image: bool = True, stdout: bool = True, save\_optimizer\_state: bool = True, validation\_epoch: int = 10, save\_checkpoint\_epoch: Optional{[}int{]} = None): Set the logging configuration.\\ API\_2: Trainer.train(gpus: str, project\_name: str): Train the model with the specified configuration. Returns a dictionary containing information about the training.\textbackslash{}n\textbackslash{}n\textbackslash{}nParameters\textbackslash{}n----------\textbackslash{}ngpus (str): GPU ids to use, separated by commas.\textbackslash{}nproject\_name (str): Project name to save the experiment.\\ API\_3: Trainer.set\_training\_config(optimizer, scheduler, epochs: int = 3, batch\_size: int = 8): Set the training configuration.\\ \\ Judgment: U\\ \\ You may need information of methods like 'set\_environment\_config()' to resolve the question.\\ But there's no relevant information about the methods. You are unanswerble to the query, so your output should be U. \\ \hline
\end{tabularx}
\caption{3-way-1shot example for NetsPressoEval}
\label{tab:2shot_prompt_netspressoeval}
\end{table*}